\documentclass[11pt,a4paper]{article}

\usepackage[utf8]{inputenc}

\usepackage[T1]{fontenc}

\usepackage[polish]{babel}

\usepackage{lmodern}

\usepackage{geometry}

\geometry{a4paper,margin=27mm}

\usepackage{microtype}

\usepackage[hidelinks]{hyperref}

% Portfolio: MSC-AR-P25; source: prospective ARC-XR research scope.

\title{Behawioralny koszt wstrzymania adaptacji XR wobec adaptacji statycznej i reaktywnej}

\author{Mateusz Pomianek}

\date{Wersja 0.3.0, 21 listopada 2026}

\begin{document}

\maketitle

\begin{abstract}

Przedmiotem opracowania jest Behawioralny koszt wstrzymania adaptacji XR wobec adaptacji statycznej i reaktywnej. Czy czasowe wstrzymanie adaptacji poprawia wynik zadania względem adaptacji statycznej i reaktywnej? Przewidziano trzy polityki sterowania HUD: statyczną, reaktywną i selektywną z jawnym wstrzymywaniem adaptacji. Głównym punktem końcowym jest odsetek błędów zadań oraz liczba niepożądanych przerwań użytkownika. Opis ma charakter wstępny i nie przedstawia niezweryfikowanych wyników.

\end{abstract}

\section{Wprowadzenie}

Rozwój adaptacyjnych systemów XR stwarza możliwość badania interakcji człowieka z informacją przestrzenną, ale samo zastosowanie AR, VR lub modeli sztucznej inteligencji nie przesądza o przewadze poznawczej czy ergonomicznej. W rozpatrywanym zagadnieniu, Behawioralny koszt wstrzymania adaptacji XR wobec adaptacji statycznej i reaktywnej, kluczowe jest rozdzielenie właściwości systemu od zmienności zachowania użytkownika. Porównanie należy zaprojektować tak, aby mierzyć jednoznaczne zjawisko, a nie tylko ogólną ocenę technologii.

\section{Cel pracy i pytanie badawcze}

Czy czasowe wstrzymanie adaptacji poprawia wynik zadania względem adaptacji statycznej i reaktywnej? Projekt wymaga krytycznego przeglądu aktualnego stanu badań oraz zdefiniowania nowości w stosunku do innych tematów portfolio XR.

\section{Zakres i odrębny wkład naukowy}
Behawioralny koszt oraz zysk wstrzymania interwencji adaptatora XR, odrębny od kalibracji niepewności predykatora P17. Test skutku abstencji na zachowanie człowieka, nie kolejna krzywa jakości selektywnego predyktora.
\section{Hipotezy badawcze}
\paragraph{H1.} Selektywne wstrzymanie interwencji zmniejsza liczbę błędów zadania względem reaktywnego adaptatora przy wyrównanej częstotliwości dostępu do wskazówek.
\paragraph{H2.} Wstrzymanie zmniejsza liczbę niepożądanych przerwań, ale nie zwiększa pominięć wymaganej pomocy ponad ustalony margines.
\paragraph{H3.} Zysk zachowuje się przy fałszywych sygnałach obciążenia, o ile pokrycie decyzji przekracza określone minimum.
\section{Projekt badania i plan eksperymentów}
\subsection{A. Operacjonalizacja i kwalifikacja danych}
Behawioralny koszt oraz zysk wstrzymania interwencji adaptatora XR, odrębny od kalibracji niepewności predykatora P17. Jednostka wnioskowania: uczestnik × zadanie z możliwością interwencji lub jej odmowy. Przed badaniem należy określić źródła, wersje urządzeń, czas akwizycji, etykiety i niezależną referencję. Nie liczyć zależnych okien i klatek jako niezależnych prób.
\subsection{B. Eksperyment porównawczy}
Porównać trzy polityki na tej samej scenie i przy identycznym źródle estymacji: brak adaptacji, adaptacja reaktywna i selektywna. Wynik pierwszorzędowy: różnica odsetka błędnych kroków zadania między polityką selektywnej odmowy adaptacji a silniejszą ze statycznej lub reaktywnej. Zamrozić losowanie warunków, progi, mianownik punktu końcowego i kryteria wyłączeń bez dostępu do oceny końcowej.
\subsection{C. Kontrole i analiza odporności}
Raportować odmowy i brak udzielenia wymaganej wskazówki, wraz z niezależną etykietą błędu czynności. Ustalony próg nie może zostać wybrany na podstawie wyników uczestników w teście; rozdzielić koszt pominięcia od kosztu nadmiarowej zmiany. Utrata sygnału, abstencja, przerwanie zadania i dane brakujące są odrębnymi kategoriami wyniku.
\subsection{D. Walidacja na niezależnym materiale}
Walidacja musi zachować podział po uczestniku, sesji, scenie lub zadaniu. Próbny model hierarchiczny poprawności i liczby interwencji z uczestnikiem i zadaniem; analiza ryzyka przy pokryciu. Raportować wyniki negatywne i rozbieżności z założeniami.
\section{Estymand, analiza statystyczna i decyzje}
różnica odsetka błędnych kroków zadania między polityką selektywnej odmowy adaptacji a silniejszą ze statycznej lub reaktywnej. Jednostka: uczestnik × zadanie z możliwością interwencji lub jej odmowy. Próbny model hierarchiczny poprawności i liczby interwencji z uczestnikiem i zadaniem; analiza ryzyka przy pokryciu. H1 jest planowanym kontrastem głównym; H2 i H3 pozostają pomocnicze do zatwierdzenia hierarchii testów oraz kontroli wielokrotności. Wartości marginesów, mocy, minimalnego efektu istotnego praktycznie i liczebności mają wynikać z pilotażu, a nie z post hoc wyboru wyniku.
\section{Stan badań i granica nowości}
\begin{itemize}
\item Existing adaptive physiological feedback in AR means adapting an interface is not by itself a novel contribution. \cite{xrprior1}.
\item Explanation design in adaptive interfaces is prior art; response to veridical versus spurious explanations needs separate evaluation. \cite{xrprior2}.
\end{itemize}
Test skutku abstencji na zachowanie człowieka, nie kolejna krzywa jakości selektywnego predyktora. Odrębność na tle opublikowanych metod i pozostałych artykułów wymaga osobnego audytu.
\section{Warunki bezpieczeństwa, sprzętu i zgodności}
To prospektywna koncepcja, nie zgoda na eksperyment, prerejestracja lub wynik empiryczny. Praca z uczestnikami, interwencje stresowe i pomiar biometryczny wymagają przeglądu etycznego, świadomej zgody i ochrony danych. Quest 3 w wersji standardowej nie posiada eye trackingu, a Tobii Eye Tracker 5 jest urządzeniem ekranowym; nie zastępować jednego stanowiska drugim.
\section{Materiał i metody}

Przewidziano trzy polityki sterowania HUD: statyczną, reaktywną i selektywną z jawnym wstrzymywaniem adaptacji. Wymagane jest zdefiniowanie jednostki analizy, zasad doboru przypadków, wersji scenariusza i źródeł pomiaru. Rozkład warunków oraz reguły wykluczenia prób należy określić przed pomiarem. W badaniach z udziałem ludzi niezbędna jest uprzednia właściwa ocena etyczna i świadoma zgoda. W testach wyłącznie symulacyjnych należy oddzielić prawdę symulacyjną od obserwacji empirycznej.

\subsection{Wynik główny i kryteria jakości}

Odsetek błędów zadań oraz liczba niepożądanych przerwań użytkownika. Pomiar powinien zachowywać jawne jednostki, znaczniki czasu i reguły oceny poprawności. Przy modelach predykcyjnych stosować rozłączne partycje oceny; przy eksperymentach interwencyjnych oddzielić skuteczność predykcji od przyczynowego wpływu interfejsu.

\section{Plan analizy}

Porównanie kontrastów polityk z kontrolą gęstości komunikatów, kolejności i trudności zadania. W wynikach należy raportować wielkość efektu lub błąd estymacji z przedziałem niepewności, liczbą obserwacji i odsetkiem odrzuconych prób. Analizy dodatkowe wymagają jasnego oznaczenia jako eksploracyjne. Nie wybierać progów decyzyjnych ani modelu na podstawie obserwacji ze zbioru testowego.

\section{Wyniki}

W obecnej wersji nie raportuje się wyników pomiarów, tabel statystycznych ani parametrów wyznaczonych z nieudokumentowanych danych. To kierunek badawczy, a nie zatwierdzony protokół badania ani zgłoszony artykuł.

\section{Dyskusja}

Interpretacja przyszłych wyników musi uwzględnić alternatywne mechanizmy i zależność wniosków od jakości danych, warunków zadania oraz ograniczeń sprzętu. Wartość potencjalnego rozwiązania zależy nie tylko od trafności modelu, lecz także od jego opóźnienia, stabilności i wpływu na zachowanie człowieka. Wyniki testów syntetycznych nie mogą być automatycznie uogólniane na populację uczestników.

\section{Ograniczenia badania}

Wstrzymanie może prowadzić do braku potrzebnej pomocy; należy oceniać zarówno błędy interwencji, jak i zaniechania. Weryfikacja skutku polityki odmowy w badaniu użytkowników; osobne pytanie wobec kalibracji predykatora XR-IDEA-11. Żadna część szkicu nie dowodzi przeprowadzenia zbierania danych, uzyskania zgód etycznych ani gotowości manuskryptu do zgłoszenia.

\section{Wnioski}

Przedstawiono problem, kontrast i sposób jego planowanej oceny. Rozstrzygnięcie postawionego pytania wymaga przygotowania protokołu, zweryfikowanych źródeł oraz poprawnie udokumentowanych wyników. Bez tych materiałów nie można sformułować konkluzji empirycznych.

\begin{thebibliography}{9}
\bibitem{xrprior1} Designing adaptive feedback systems for managing cognitive load in augmented reality (2025), Design Science 11, e49. \url{https://doi.org/10.1017/dsj.2025.10040}.
\bibitem{xrprior2} EXAD: Framework for Designing Explanations in Adaptive User Interfaces (2025), IEEE Conference. \url{https://ieeexplore.ieee.org/abstract/document/11017085}.
\end{thebibliography}
\end{document}

